\chapter{La segunda vía}
\label{ch:talent}

\section{La afirmación, enunciada con cuidado}

Este es el capítulo más especulativo del libro y conviene que lo anuncie antes de
argumentar nada.

La afirmación es que Cuba puede construir un segundo sector exportador, junto al turismo,
a partir de gente educada que venda servicios a través de una frontera sin salir de la
isla ---software, ingeniería, diseño, servicios administrativos, trabajo con datos,
contenido en español, servicios profesionales remotos--- y que ese sector podría
plausiblemente alcanzar el mismo orden de magnitud que el turismo en una década.

La afirmación no es que Cuba se convierta en una potencia de inteligencia artificial. No
lo será. Entrenar modelos grandes exige electricidad que Cuba no tiene, capital que no
puede reunir y equipamiento que no puede comprar legalmente. Cualquier propuesta en la que
Cuba construya centros de datos no es una propuesta sobre Cuba. Lo que Cuba sí puede hacer
plausiblemente es aportar \emph{trabajo} al mercado global de servicios que la ola
tecnológica actual ha ampliado ---la anotación, la evaluación, la localización, la
integración, las herramientas y la ingeniería de software corriente que están alrededor de
los modelos y no dentro de ellos--- y puede hacerlo en español, en el huso horario de Nueva
York, con una fuerza de trabajo que sabe leer.

Vale la pena poner la aritmética sobre la página porque es la razón de ser del capítulo.
Cien mil cubanos ---cerca del dos por ciento de la fuerza laboral--- ganando un promedio de
dieciocho mil dólares al año en servicios remotos generarían aproximadamente 1.800 millones
de dólares anuales en divisas.\unv{} Es el mismo orden que los ingresos turísticos cubanos
en un buen año, desde un sector que no requiere hoteles, ni playa, ni abasto de agua, ni
combustible, y casi nada de capital. La cifra por trabajador está muy por debajo de lo que
gana un trabajador comparable en un país de altos ingresos y muy por encima de cualquier
cosa disponible hoy en Cuba, que es precisamente por lo que el mercado se equilibraría.

\section{Por qué no ha ocurrido}

Cuba lleva treinta años teniendo trabajadores técnicos educados, subempleados y capaces en
inglés, y el mercado global de servicios remotos existe desde hace veinte. El sector no
apareció. Las razones son específicas, e identificar cuál es la vinculante es todo el
contenido de este capítulo.

No es el talento. La educación técnica cubana es el activo que describió el
capítulo~\ref{ch:socialstate}, y los programadores cubanos que trabajan en el exterior son
competitivos.

No es, en lo principal, el embargo. Los clientes norteamericanos están restringidos, pero
España, México, Argentina, Canadá, Alemania y todo el mundo hispanohablante no lo están, y
ninguno de ellos compra servicios cubanos a escala tampoco.

No es, principalmente, la conectividad ---aunque la conectividad es una restricción real y
la sección siguiente trata de ella---. Hay muchos países que exportan servicios con menos
ancho de banda del que La Habana tiene desde 2018.

\emph{La restricción vinculante es que un cubano no puede contratar de forma legal y
fiable con un cliente extranjero y cobrar en una cuenta que controle.} Hasta 2021 un
operador privado cubano no tenía personalidad jurídica y por tanto no podía firmar un
contrato como entidad. La lista de actividades permitidas estaba escrita en el vocabulario
de los oficios y no tenía categoría para un desarrollador de software ni para un estudio de
diseño. Las cuentas en divisa estaban restringidas. Los rieles internacionales de tarjetas
y transferencias no llegaban a los individuos. Y a los profesionales formados a costa del
Estado se les prohibió durante buena parte del período ejercer su profesión de forma
privada.

El resultado es que los cubanos que hacen este trabajo lo hacen informalmente, a través de
la cuenta extranjera de un pariente, o a través de un intermediario en Miami o Madrid que se
lleva una comisión y se queda con la relación con el cliente ---o, cada vez más, emigran, lo
que convierte un sector exportador potencial en una pérdida permanente---.

Esta es la buena noticia del capítulo. La restricción vinculante es jurídica, es doméstica,
no le cuesta nada al Estado retirarla, y no requiere el permiso de nadie.

\begin{design}{El derecho a vender el propio trabajo}
\label{d:contract}
\begin{rules}
\rl{1}{El derecho central} Todo residente puede contratar directamente con un cliente
extranjero, en cualquier moneda, sin licencia, intermediario ni autorización previa, y puede
recibir el pago en una cuenta a su nombre y gastar desde ella. Es la cláusula individual más
importante de la estrategia de crecimiento de la Parte~III y es una derogación, no un
programa.

\rl{2}{Profesiones incluidas} Se extiende explícitamente a las profesiones formadas a costa
del Estado: ingenieros, médicos, arquitectos, juristas, científicos. La exclusión de los
profesionales del trabajo por cuenta propia desde 1993 es la razón de que el sector privado
cubano sea un conjunto de oficios, y la razón de que sus trabajadores más valiosos tuvieran
solo dos opciones: el Estado o el aeropuerto.

\rl{3}{Ninguna lista enumerada} Según el Diseño~\ref{d:commit}.2, la exportación de
servicios se permite por derecho. Cualquier lista de actividades permitidas en este dominio
quedaría obsoleta antes de publicarse, lo que es un argumento general contra las listas y uno
decisivo aquí.

\rl{4}{Un impuesto sencillo que valga la pena pagar} Los ingresos por servicios
transfronterizos pagan la tasa plana baja del Diseño~\ref{d:tax}.6, con declaración anual,
sin deducciones y sin contribución social aparte. Una base móvil gravada con dureza no se
grava; una base móvil gravada de forma ligera y sencilla se registra, se cuenta, y con el
tiempo se convierte en una clientela del sistema que la registró.

\rl{5}{Equipamiento} Computadoras, teléfonos, equipos de red y sus piezas entran libres de
arancel y de licencia, según el Diseño~\ref{d:genrights}.6. Cuba tiene hoy ciudadanos que
compran sus propias herramientas con dinero de remesas y un régimen aduanero que lo trata
como contrabando.

\rl{6}{Clasificación} Los servicios transfronterizos se clasifican y contabilizan como
exportaciones en las cuentas nacionales y en la balanza de pagos, lo que suena a trámite y no
lo es: lo que no se mide no se defiende en el presupuesto, y este sector va a necesitar
defensa frente a los ministerios de cuyo personal se nutre.
\end{rules}
\end{design}

\section{La conectividad}

Cuba se conectó tarde y caro. Los datos móviles llegaron en diciembre de 2018; la capacidad
internacional ha corrido principalmente por un único cable submarino tendido desde Venezuela
en 2011, que es a la vez una restricción de ancho de banda y un punto único de fallo; la
empresa estatal de telecomunicaciones es un monopolio; y el precio de los datos como
proporción del ingreso cubano ha sido de los más altos del mundo, que es de lo que trataban
las protestas por el cambio de tarifas de 2025.

Ninguna segunda vía existe sobre esa base. Un cliente extranjero que contrata con un
proveedor cubano está comprando fiabilidad, y la fiabilidad es exactamente lo que no dan un
solo cable, un operador monopólico, una red eléctrica frágil y un Estado que cortó los datos
móviles a escala nacional el 11 de julio de 2021.

\begin{design}{La conectividad como infraestructura}
\label{d:connect}
\begin{rules}
\rl{1}{Redundancia} Al menos dos rutas internacionales adicionales, por trazados físicos
diversos y con puntos de amarre diversos, contratadas en condiciones competitivas
publicadas. Un cable no es una red.

\rl{2}{Acceso mayorista abierto} Cualquiera que sea la propiedad de la capacidad
internacional y troncal, el acceso a ella se vende a cualquier proveedor minorista con
licencia en condiciones publicadas y no discriminatorias, reguladas por una autoridad
independiente al modo del Diseño~\ref{d:genrights}.4. La competencia minorista es lo que mueve
el precio; la privatización completa del operador establecido no es necesaria y no debe
intentarse temprano.

\rl{3}{Ninguna interrupción política} Según el Diseño~\ref{d:minimum}.5. Se repite aquí
porque es una cláusula comercial tanto como una de libertades civiles: un país cuyo gobierno
ha demostrado que apagará internet durante un episodio político no puede vender acuerdos de
nivel de servicio, y todo cliente potencial lo sabe.

\rl{4}{Transparencia de precios} Tarifas mayoristas y minoristas publicadas, y una métrica
publicada de asequibilidad ---costo de los datos como proporción del ingreso mediano---
informada anualmente.

\rl{5}{Electricidad} Los sitios de telecomunicaciones son cargas críticas según el
Diseño~\ref{d:energyseq}.1 y están entre los primeros en recibir fotovoltaica y
almacenamiento. Una red que se cae con la red eléctrica tampoco es una red.
\end{rules}
\end{design}

\section{El régimen de residencia}

Solo después de los dos diseños anteriores tiene sentido una política de visados, y el orden
importa porque la discusión internacional sobre este asunto lo tiene al revés.

El visado de nómada digital es un instrumento modesto. La lectura honesta de la experiencia
internacional ---Estonia, Portugal, Barbados, Croacia, Costa Rica--- es que estos esquemas
atraen a unos pocos miles de personas, generan un gasto de consumo útil pero pequeño, y no
construyen por sí mismos una industria. El Welcome Stamp de Barbados fue una respuesta
competente a un desplome turístico pandémico, no una transformación económica. La
e-Residency de Estonia, el más citado de todos, ha inscrito a un gran número de personas y ha
producido un número mucho menor de empresas reales que paguen impuestos reales en Estonia.

Lo que de verdad construyó las exportaciones de servicios en los países comparadores no
fueron los visados para extranjeros. Fue que \emph{a los residentes se les permitiera y se
les hiciera posible vender su trabajo al exterior}, más un derecho previsible y un sistema
de pagos ---que es el Diseño~\ref{d:contract}---. El sector de software de Uruguay, el caso
latinoamericano más pertinente, se construyó con empresas e ingenieros nacionales y no con
nómadas entrantes.

De modo que el régimen de residencia se propone aquí por tres razones más estrechas y más
defendibles: trae gente que forma y contrata localmente; trae demanda del alojamiento de
larga estancia que el capítulo~\ref{ch:tourism} quiere hacer crecer en vez de habitaciones de
hotel; y ---según el Diseño~\ref{d:pension}.3--- más residentes en edad de trabajar es el
instrumento más barato disponible para una razón de dependencia.

\begin{design}{Residencia para trabajadores móviles}
\label{d:residence}
\begin{rules}
\rl{1}{El permiso} Un permiso de residencia renovable para quien tenga ingresos extranjeros
por trabajo remoto o por un negocio en el exterior, otorgado según criterios objetivos
publicados ---umbral de ingresos, cobertura de salud, antecedentes limpios--- con un plazo
legal de decisión y otorgamiento tácito si el plazo transcurre. Sin discrecionalidad, sin
entrevista, sin patrocinador.

\rl{2}{Certeza fiscal} Una tasa plana publicada sobre los ingresos remitidos localmente
durante un número fijo de años, legislada en vez de administrada, con el tratamiento de los
ingresos extranjeros declarado por adelantado. La incertidumbre disuade a esta población
mucho más que las tasas.

\rl{3}{Derechos ordinarios} El derecho a alquilar, a abrir una cuenta, a importar equipo
personal, a traer familia, y a usar el sistema de salud en las mismas condiciones que los
residentes pagando los mismos impuestos. Un permiso que no permite una cuenta bancaria no es
un permiso.

\rl{4}{Una vía para quedarse} Un camino definido a la residencia permanente y, para quien lo
quiera, a la naturalización. El objetivo son residentes, no visitantes; un esquema que se
reinicia cada año selecciona a gente que no echa raíces.

\rl{5}{Encadenamiento local} Una tasa reducida para los titulares del permiso que empleen o formen
formalmente a residentes cubanos, verificado de forma sencilla y auditado. Es lo que convierte
un flujo entrante en una industria doméstica, y es la diferencia entre esta propuesta y una
playa con mejor wifi.
\end{rules}
\end{design}

\section{La diáspora, y no gravar la puerta}

Dos millones de personas en el exterior, la mayoría a un vuelo de distancia, muchas con
capital y posición profesional, y un número sustancial que trabajaría con Cuba si trabajar
con Cuba fuera posible. Son a la vez los clientes iniciales naturales de este sector, los
inversionistas iniciales naturales del capítulo~\ref{ch:capital}, y el mayor depósito
individual de la gente en edad de trabajar que el capítulo~\ref{ch:premises} dice que el país
ha perdido.

\begin{design}{Circulación, no retención}
\label{d:diaspora}
\begin{rules}
\rl{1}{Ningún impuesto de salida} Cuba no grava, multa ni impone un canon de formación a los
emigrantes. El instrumento es popular entre los gobiernos que pierden trabajadores
calificados, no recauda casi nada, se elude trivialmente con no regresar, y convierte una
partida temporal en una permanente. Es además, en el caso cubano específicamente, un impuesto
sobre el agravio que este libro trata de saldar.

\rl{2}{Retorno sin penalización} Restauración plena de la residencia, los derechos de
propiedad, el registro profesional y el derecho a pensión de los cubanos que regresen,
automática y sin revisión discrecional, incluidos quienes se fueron de forma irregular o
abandonaron una misión.

\rl{3}{Doble nacionalidad} Reconocida, con derecho a tener propiedad, a invertir, a trabajar,
a votar en elecciones locales tras un requisito de residencia, y a heredar. El
capítulo~\ref{ch:capital} depende de esto; también la portabilidad de pensiones del
Diseño~\ref{d:pension}.4.

\rl{4}{Contribución voluntaria} Los cubanos en el exterior pueden contribuir voluntariamente
al sistema cubano de seguridad social y acumular derechos, según el
Diseño~\ref{d:pension}.4. Esto convierte una diáspora en una base de cotizantes, cosa que
ningún otro instrumento hace.

\rl{5}{Profesores y clientes} Programas para traer a académicos y profesionales de la
diáspora en estancias cortas de docencia, clínica y asesoría, sin prueba política,
administrados por las universidades y no por los ministerios.
\end{rules}
\end{design}

\section{Las universidades}

El sector que aquí se propone necesita una cantera, y el capítulo~\ref{ch:socialstate}
identificó la brecha específica: Cuba produce ingenieros, médicos y matemáticos excelentes y
muy poca gente formada en las disciplinas de las que en realidad viven una transición
institucional y una economía de servicios ---derecho mercantil, contabilidad, finanzas,
administración, economía aplicada y ciencias sociales empíricas---.

\begin{design}{Capacidad académica}
\label{d:university}
\begin{rules}
\rl{1}{Autonomía} Las universidades nombran a su propio personal, fijan sus propios planes de
estudio y admiten a sus propios estudiantes según criterios académicos publicados. El
expediente político no es criterio de admisión ni de nombramiento, y la prohibición es
justiciable.

\rl{2}{Las disciplinas ausentes} Un programa financiado para construir facultades de derecho,
contabilidad, finanzas, administración y economía aplicada, con asociación extranjera,
nombramientos visitantes desde la diáspora, y acreditación profesional reconocida
internacionalmente. Es una restricción de una década sobre todo lo de la Parte~III ---el
Diseño~\ref{d:courts}.6 necesita jueces mercantiles, y el capítulo~\ref{ch:welfare} necesita
administradores tributarios--- y obliga desde el año uno.

\rl{3}{Asociación extranjera} Programas conjuntos, títulos dobles y sedes en asociación con
universidades extranjeras, permitidos por derecho y regulados solo por calidad.

\rl{4}{Publicar y ser leídos} Trabajo académico cubano publicado abiertamente, y académicos
cubanos libres de publicar en el exterior, colaborar y viajar. El desmantelamiento del Centro
de Estudios sobre América en 1996 le enseñó a las ciencias sociales cubanas lo que le ocurre a
la investigación que analiza la política pública, y esa lección tiene que desaprenderse
públicamente antes de que nadie haga el trabajo que necesitan las fases posteriores de este
libro.
\end{rules}
\end{design}

\section{Lo que este capítulo no puede prometer}

Tres límites honestos.

El sector es pequeño hasta que se arregle la electricidad. Todo lo de aquí está aguas abajo
del capítulo~\ref{ch:energy}, y un trabajador remoto en una provincia con dieciséis horas de
apagón no puede atender una llamada.

El sector compite exactamente por la gente que el Estado más necesita. Cada ingeniero que
pasa al trabajo remoto por dieciocho mil dólares al año es un ingeniero que no está en el
ministerio, la empresa eléctrica o la universidad con un salario estatal ---por lo cual el
Diseño~\ref{d:pay} es precondición de este capítulo y no un tema aparte---.

Y el sector es un rival de la emigración, no un complemento. Su justificación real es que le
ofrece a una persona el ingreso por el que se iría sin exigirle que se vaya. Si funciona, el
flujo de salida se frena; si el flujo de salida no se frena, no está funcionando, y el
capítulo~\ref{ch:sequence} convierte la migración neta en uno de los indicadores por los que
puede juzgarse que este libro está fracasando.
